% =====================================================================
% Experiment 4 (Vlad) — Compute-Matched Efficiency, SDXL x Pick-a-Pic
% Paste-ready for the group report (ACM sigconf). Numbers trace to
% report/tables/vlad_efficiency.md and outputs/vlad_exp8_sdxl_pickapic/.
% Figures: copy these PNGs into the Overleaf figures/ directory:
%   report/figures/vlad_exp8_sdxl_pickapic/exp8_hps_vs_time.png   -> figures/exp8_hps_vs_time.png
%   report/figures/vlad_exp8_sdxl_pickapic/exp8_qual_standard.png -> figures/exp8_qual_standard.png
%   report/figures/vlad_exp8_sdxl_pickapic/exp8_qual_zigzag.png   -> figures/exp8_qual_zigzag.png
% Use the same heading macro as the other experiments (shown as \subsection*).
% If the 8-column table overflows the single column, change {table} -> {table*}.
% =====================================================================

\subsection*{Experiment 4 --- Compute-Matched Efficiency: SDXL on Pick-a-Pic (Vlad)}

Z-Sampling adds an inversion and re-denoising pass at each step, so a fair comparison must
fix total computation, not the number of sampling steps. We reproduce the paper's
time-matched efficiency study (Figure~10) on Stable Diffusion XL base 1.0 over the first 100
Pick-a-Pic prompts (seed 42). For each standard budget $T_s \in \{20, 30, 40, 50\}$ we
compare standard DDIM at $T_s$ steps against Z-Sampling with $T_z = T_s/2$ and
$\lambda = T_s/4$, which equalises the noise-prediction count ($T_z + 2\lambda = T_s$). Both
arms share the initial latent, scheduler (DDIM with inversion), guidance ($\gamma_1 = 5.5$,
$\gamma_2 = 0$) and resolution ($1024^2$). We score all 400 pairs with HPS\,v2.1 and record
per-image time on Apple silicon (MPS); a few cells whose wall-clock was corrupted when the
laptop was suspended mid-run (timings above twice the per-budget median) are excluded from
the time summary, while every image is kept for scoring.

At each budget the two arms take almost identical median time (for example, 72.2\,s at
$T_s = 40$), confirming the match, yet Z-Sampling scores higher at every budget and is
preferred on 81\% to 91\% of pairs (85.2\% overall; Table~\ref{tab:vlad-efficiency}). The
pooled delta is $+0.0148$ with a 95\% bootstrap confidence interval of $[+0.0131, +0.0165]$,
strictly positive at every budget and larger than the $+0.0110$ we measure for Stable
Diffusion 2.1. Read across budgets, this recovers the paper's Pareto result
(Figure~\ref{fig:vlad-efficiency}): the best standard score, $0.2727$, needs $T_s = 50$
($\sim$151\,s), whereas Z-Sampling reaches $0.2779$ at $T_s = 30$ ($\sim$53\,s), matching the
best standard quality at roughly one third of the time. This reproduces the direction of the
paper's ``36\% less computational time'' claim, though the absolute saving is likely
optimistic, as the $T_s = 50$ reference budget shows signs of thermal throttling on this
device.

Two caveats temper this. First, op-count parity is only an approximation of true cost, since
the interleaved inversion changes the scheduler and attention access pattern and the gap
could widen on other stacks. Second, HPS\,v2.1 is the sole judge, and because Z-Sampling
amplifies exactly the prompt-alignment signal preference models reward, a single-metric win
may partly reflect that alignment; corroboration with PickScore, ImageReward, or AES lay
beyond our compute budget, though other experiments in this report evaluate SDXL under those
metrics. The core finding still holds: winning at \emph{equal} compute indicates that
spending the budget on $\gamma_1\!\to\!\gamma_2\!\to\!\gamma_1$ self-reflection over the
early, layout-setting timesteps beats additional plain denoising
(Figure~\ref{fig:vlad-qualitative}). The single seed, fixed $\lambda$, and absent resampling
arm are deliberate scope limits under constrained hardware.

\begin{table}[t]
\centering
\caption{Compute-matched efficiency comparison on Stable Diffusion XL with Pick-a-Pic
(100 prompts per budget, seed 42). Z-Sampling uses $T_z = T_s/2$ and $\lambda = T_s/4$ so
that the two methods perform an equal number of noise predictions. Time is the median
seconds per image on MPS, HPS is the mean HPS\,v2.1 score, and the last column reports the
fraction of pairs on which Z-Sampling scores higher.}
\label{tab:vlad-efficiency}
\small
\begin{tabular}{rrrrrrrr}
\toprule
$T_s$ & $T_z$ & $\lambda$ & Std s/img & Z s/img & Std HPS & Z HPS & Z win \\
\midrule
20 & 10 & 5  & 36.8  & 36.9  & 0.2574 & 0.2720 & 81\% \\
30 & 15 & 7  & 54.4  & 52.7  & 0.2645 & 0.2779 & 82\% \\
40 & 20 & 10 & 72.2  & 72.2  & 0.2712 & 0.2860 & 87\% \\
50 & 25 & 12 & 151.4 & 149.0 & 0.2727 & 0.2891 & 91\% \\
\bottomrule
\end{tabular}
\end{table}

\begin{figure}[t]
\centering
\includegraphics[width=\columnwidth]{figures/exp8_hps_vs_time.png}
\caption{Mean HPS\,v2.1 against median generation time per image (reproduction of
Figure~10 of the original paper), SDXL on Pick-a-Pic. At every compute-matched budget the
Z-Sampling curve lies above standard sampling, and at $T_s = 30$ it exceeds the best
standard point. The error bars are per-arm 95\% confidence intervals and reflect score
spread across prompts (dominated by prompt difficulty, hence the overlap); statistical
significance is established on the paired per-prompt delta of
Table~\ref{tab:vlad-efficiency}, whose bootstrap interval excludes zero at every budget.}
\label{fig:vlad-efficiency}
\end{figure}

\begin{figure}[t]
\centering
\begin{minipage}{0.49\columnwidth}\centering
\includegraphics[width=\linewidth]{figures/exp8_qual_standard.png}\\
{\small Standard, $T_s{=}40$}
\end{minipage}\hfill
\begin{minipage}{0.49\columnwidth}\centering
\includegraphics[width=\linewidth]{figures/exp8_qual_zigzag.png}\\
{\small Z-Sampling, $T_z{=}20$, $\lambda{=}10$}
\end{minipage}
\caption{A representative compute-matched pair (prompt: ``An elephant in a room full of
glassware'', seed 42, budget $T_s = 40$; HPS delta $+0.014$, at the per-budget mean). At
equal compute Z-Sampling (right) renders more coherent detail and composition than standard
sampling (left). This is a typical case selected near the mean delta, not the largest gain.}
\label{fig:vlad-qualitative}
\end{figure}

% =====================================================================
% Appendix B — Vlad's work-division row (uncomment and paste into Table 5)
% =====================================================================
% Vlad & Implemented the compute-matched efficiency runner and the aggregation pipeline
% (\texttt{run\_vlad\_exp8\_efficiency.py}, \texttt{summarize\_vlad\_efficiency.py}).
% Conducted Experiment~4 (standard sampling versus compute-matched Z-Sampling on Stable
% Diffusion XL with Pick-a-Pic, evaluated with HPS\,v2.1 and runtime; reproduction of
% Figure~10). Presented the opening section (paper motivation, models, and datasets) at
% the end-term presentation. Wrote the Experiment~4 section of the report. \\
